\documentclass[notes=show]{beamer}
%%%%%%%%%%%%%%%%%%%%%%%%%%%%%%%%%%%%%%%%%%%%%%%%%%%%%%%%%%%%%%%%%%%%%%%%%%%%%%%%%%%%%%%%%%%%%%%%%%%%%%%%%%%%%%%%%%%%%%%%%%%%%%%%%%%%%%%%%%%%%%%%%%%%%%%%%%%%%%%%%%%%%%%%%%%%%%%%%%%%%%%%%%%%%%%%%%%%%%%%%%%%%%%%%%%%%%%%%%%%%%%%%%%%%%%%%%%%%%%%%%%%%%%%%%%%
\AtBeginSubsection[
  {\frame<beamer>{\frametitle{Outline}   
    \tableofcontents[currentsection,currentsubsection]}}%
]%
{
  \frame<beamer>{ 
    \frametitle{Outline}   
    \tableofcontents[currentsection,currentsubsection]}
}
\usepackage{ulem}
\usepackage{pdflscape}
\setbeamertemplate{caption}{\insertcaption}
\usepackage{color,xcolor,ucs}
\usepackage{beamerthemesplit}
\usepackage{graphicx}
\usepackage{amsmath}
\usepackage{adjustbox}
\usepackage{pdfpages} 
\definecolor{mintgreen}{rgb}{0.6, 1.0, 0.6}
\makeatletter
\@addtoreset{subfigure}{framenumber}% subfigure counter resets every frame
\makeatother
\newcommand*{\Scale}[2][4]{\scalebox{#1}{$#2$}}%
\newcommand*{\Resize}[2]{\resizebox{#1}{!}{$#2$}}%
\usepackage{amsfonts}
\usepackage{booktabs}
\usepackage{color,soul}
\usepackage{pdfpages}
\setboolean{@twoside}{false}
\usepackage{amsmath}
\usepackage{pdflscape} 
\usepackage{eurosym}
\usepackage{amstext} % for \text
\DeclareRobustCommand{\officialeuro}{%
  \ifmmode\expandafter\text\fi
  {\fontencoding{U}\fontfamily{eurosym}\selectfont e}}
\usepackage[caption = false]{subfig}
\usepackage{appendixnumberbeamer} 
\usepackage{graphicx}
\usepackage{mathpazo}
\usepackage{hyperref}
\usepackage{multimedia}
\usepackage{graphicx}
\usepackage{multirow}
\usepackage{subfig}
\usepackage{verbatim}
\usepackage{booktabs, multicol, multirow}
\setcounter{MaxMatrixCols}{10}
\input{Macros.tex}
\AtBeginSection[]
{
  \begin{frame}[noframenumbering]
    \frametitle{Outline for section \thesection}
    \tableofcontents[currentsection]
  \end{frame}
}


\usepackage{xcolor,soul}
\renewcommand<>{\hl}[1]{\only#2{\beameroriginal{\hl}}{#1}}

% https://tex.stackexchange.com/questions/41683/why-is-it-that-coloring-in-soul-in-beamer-is-not-visible
\makeatletter
\newcommand\SoulColor{%
  \let\set@color\beamerorig@set@color
  \let\reset@color\beamerorig@reset@color}
\makeatother
\SoulColor
\AtBeginSection{\frame{\sectionpage}}

\newsavebox\mybox
\newenvironment{aquote}[1]
  {\savebox\mybox{#1}\begin{quote}\openautoquote\hspace*{-.7ex}}
  {\unskip\closeautoquote\vspace*{1mm}\signed{\usebox\mybox}\end{quote}}

\newtheorem{proposition}[theorem]{Proposition}
\newtheorem{Assumption}[theorem]{Assumption}
\newenvironment{stepenumerate}{\begin{enumerate}[<+->]}{\end{enumerate}}
\newenvironment{stepitemize}{\begin{itemize}[<+->]}{\end{itemize} }
\newenvironment{stepenumeratewithalert}{\begin{enumerate}[<+-| alert@+>]}{\end{enumerate}}
\newenvironment{stepitemizewithalert}{\begin{itemize}[<+-| alert@+>]}{\end{itemize} }
\usetheme{Madrid}
\makeatletter
\setbeamertemplate{footline}
{
  \leavevmode%
  \hbox{%
  \begin{beamercolorbox}[wd=.333333\paperwidth,ht=2.25ex,dp=1ex,center]{author in head/foot}%
    \usebeamerfont{author in head/foot}\insertshortauthor~~\beamer@ifempty{\insertshortinstitute}{}{(\insertshortinstitute)}
  \end{beamercolorbox}%
  \begin{beamercolorbox}[wd=.333333\paperwidth,ht=2.25ex,dp=1ex,center]{title in head/foot}%
    \usebeamerfont{title in head/foot}\insertshorttitle
  \end{beamercolorbox}%
  \begin{beamercolorbox}[wd=.333333\paperwidth,ht=2.25ex,dp=1ex,right]{date in head/foot}%
    \usebeamerfont{date in head/foot}\insertshortdate{}\hspace*{2em}
%    \insertframenumber{} / \inserttotalframenumber\hspace*{2ex} % DELETED
  \end{beamercolorbox}}%
  \vskip0pt%
}
\makeatother
\usepackage{pdfpages}
\begin{document}
	\title[]{\large Regression Discontinuity
	\vspace{0.5cm}
	}
	\author[Raffaele Saggio - Econ 628]{Raffaele Saggio}
	\institute[]{UBC}
	\date{\today \\
}
	\maketitle


\begin{frame}{Introduction}

\begin{itemize}
\item The \textbf{regression discontinuity design }(RD) is one of the most
popular quasi-experimental research designs in contemporary applied
research\medskip{}
\item RD methods can be applied when a researcher has specific information
about the rules determining the treatment of interest\medskip{}
\item This lecture describes the RD framework and implementation\medskip{}
\item See MHE; Imbens and Lemieux (2008), Lee and Lemieux (2010), and Imbens
and Wooldridge (2009) for recent overviews of RD
\end{itemize}
\end{frame}


\begin{frame}{Sharp RD: Basic Setup}

\begin{itemize}
\item Consider a setting with a treatment $T_{i}\in\{0,1\}$, and potential
outcomes $Y_{i}(1)$ and $Y_{i}(0)$\medskip{}
\item Suppose the treatment is a \emph{deterministic} function of an observed
covariate $R_{i}$, such that
\end{itemize}
\begin{center}
$T_{i}=1\left\{ R_{i}>c\right\} $. 
\par\end{center}
\begin{itemize}
\item $R_{i}$ is called the \textbf{running variable}\medskip{}
\item This is a \textbf{sharp} \textbf{RD} because the probability of treatment
switches from zero to one at the threshold\medskip{}
\item Example: Scholarship awarded to students who score above a test score
threshold (Thistlethwaite and Campbell 1960)
\end{itemize}
\end{frame}

\begin{frame}{Sharp RD: Basic Setup}

\begin{itemize}
\item We get to observe $Y_{i}(1)$ when $R_{i}>c$ and $Y_{i}(0)$ when
$R_{i}\leq c$\medskip{}
\item Basic idea of RD design: Compare observations just above and just
below the threshold to infer treatment effect\medskip{}
\item Intuitively, the treatment may be as good as randomly assigned for
individuals in the neighborhood of $R_{i}=c$, so comparing treated
and nontreated reveals a treatment effect
\end{itemize}
\end{frame}

\setbeamercolor{background canvas}{bg=} \includepdf{aux/RDD.pdf}
\begin{frame}{Sharp RD Identification}

\begin{itemize}
\item Key RD identifying assumption: potential outcomes are smooth at the
threshold\medskip{}
\item Formally:
\end{itemize}
\begin{center}
${\displaystyle \lim_{r\rightarrow c^{+}}E\left[Y_{i}(t)|R_{i}=r\right]}={\displaystyle \lim_{r\rightarrow c^{-}}E\left[Y_{i}(t)|R_{i}=r\right]},\ t\in\{0,1\}$
\par\end{center}
\begin{itemize}
\item Potential outcome CEFs must be continuous at the threshold\medskip{}
\item The population just below must not be discretely different from the
population just above
\medskip{}
\item The only thing that is allowed to change disproportionately at the threshold is the assignment probability of treatment 
\end{itemize}
\end{frame}

\begin{frame}{Sharp RD Identification}

\begin{itemize}
\item If this assumption holds we have
\end{itemize}
\begin{center}
${\displaystyle \lim_{r\rightarrow c^{+}}E\left[Y_{i}|R_{i}=r\right]-{\displaystyle \lim_{r\rightarrow c^{-}}E\left[Y_{i}|R_{i}=r\right]}}$
\par\end{center}

\begin{center}
$={\displaystyle \lim_{r\rightarrow c^{+}}E\left[Y_{i}(1)|R_{i}=r\right]-{\displaystyle \lim_{r\rightarrow c^{-}}E\left[Y_{i}(0)|R_{i}=r\right]}}$
\par\end{center}

\begin{center}
$={\displaystyle E\left[Y_{i}(1)|R_{i}=c\right]-{\displaystyle E\left[Y_{i}(0)|R_{i}=c\right]}}$
\par\end{center}

\begin{center}
$=E\left[Y_{i}(1)-Y_{i}(0)|R_{i}=c\right]$
\par\end{center}
\begin{itemize}
\item When potential outcomes are smooth around the threshold, a comparison
of individuals just above and just below yields the average treatment
effect for individuals at the threshold\medskip{}
\item This identification argument is nonparametric: we didn't have to assume
anything about the distribution of $Y_{i}(t)$ other than continuity
of CEFs
\end{itemize}
\end{frame}

\begin{frame}{Sharp RD Interpretation}

\begin{itemize}
\item Note that in a sharp RD design we have the CIA condition
\end{itemize}
\begin{center}
$(Y_{i}(1),Y_{i}(0))\bot T_{i}|R_{i}$
\par\end{center}
\begin{itemize}
\item CIA holds trivially because there is no variation in $T_{i}$ conditional
on $R_{i}$\medskip{}
\item But there is also no common support -- $T_{i}$ is either always
zero or always one conditional on $R_{i}$\medskip{}
\item RD estimation is a local extrapolation outside the support of the
data to predict the mean treated and untreated potential outcomes
at $R_{i}=c$
\end{itemize}
\end{frame}

\begin{frame}{RD Estimation}

\begin{itemize}
\item Implementing a sharp RD requires estimating the right- and left-hand
limits
\end{itemize}
\begin{center}
${\displaystyle \lim_{r\rightarrow c^{+}}E\left[Y_{i}|R_{i}=r\right]}$
\par\end{center}

\begin{center}
${\displaystyle \lim_{r\rightarrow c^{-}}E\left[Y_{i}|R_{i}=r\right]}$
\par\end{center}
\begin{itemize}
\item Since extrapolation is central to the RD design, functional form used
to approximate $E\left[Y_{i}|R_{i}\right]$ really matters\medskip{}
\item Not enough to just do OLS and rely on approximation theorems\medskip{}
\item Insufficiently flexible specification of the CEF runs the risk of
mistaking nonlinearity for treatment effect\medskip{}
\item But overly flexible specification reduces precision and runs the risk
of overfitting
\end{itemize}
\end{frame}

\setbeamercolor{background canvas}{bg=} \includepdf{aux/rdmhe_1.pdf}

\setbeamercolor{background canvas}{bg=} \includepdf{aux/rdmhe_2.pdf}

\setbeamercolor{background canvas}{bg=} \includepdf{aux/rdmhe_3.pdf}
\begin{frame}{RD Estimation}

\begin{itemize}
\item Two general approaches to estimation:\medskip{}

\begin{itemize}
\item Global parametric\medskip{}
\item Local nonparametric
\end{itemize}
\end{itemize}
\end{frame}

\begin{frame}{Global Polynomial}

\begin{itemize}
\item The global approach uses parametric functions to approximate $E\left[Y_{i}|R_{i}\right]$,
typically polynomials\medskip{}
\item Estimating equation:
\end{itemize}
\begin{center}
$Y_{i}=\alpha+\beta1\left\{ R_{i}>c\right\} +\sum_{k=1}^{K}\gamma_{0k}\left(R_{i}-c\right)^{k}+\sum_{k=1}^{K}\gamma_{1k}1\left\{ R_{i}<c\right\} \left(R_{i}-c\right)^{k}+\epsilon_{i}$
\par\end{center}
\begin{itemize}
\item This specification uses a $K$-th order polynomial with coefficients
that differ on each side of the threshold\medskip{}
\item The parameter $\beta$ measures the jump at the threshold
\end{itemize}
\end{frame}

\begin{frame}{Local Linear Regression}

\begin{itemize}
\item More modern approach: use nonparametric techniques to approximate
the left- and right-hand limits\medskip{}
\item Local linear regression:
\end{itemize}
\begin{center}
$\left(\hat{\alpha}_{0},\hat{\delta}_{0}\right)={\displaystyle \arg\min_{\alpha_{0},\delta_{0}}}{\displaystyle \sum_{i}}1\left\{ R_{i}\leq c\right\} K\left(\tfrac{R_{i}-c}{h}\right)\left(Y_{i}-\alpha_{0}-\delta_{0}\left(R_{i}-c\right)\right)^{2}$
\par\end{center}
\begin{itemize}
\item Here $K(\cdot)$ is a symmetric kernel function maximized at 0, and
$h$ is a bandwidth that vanishes asymptotically
\end{itemize}
\end{frame}

\begin{frame}{Local Linear Regression}

\begin{center}
$\left(\hat{\alpha}_{0},\hat{\delta}_{0}\right)={\displaystyle \arg\min_{\alpha_{0},\delta_{0}}}{\displaystyle \sum_{i}}1\left\{ R_{i}\leq c\right\} K\left(\tfrac{R_{i}-c}{h}\right)\left(Y_{i}-\alpha_{0}-\delta_{0}\left(R_{i}-c\right)\right)^{2}$
\par\end{center}
\begin{itemize}
\item This is weighted least squares using observations to the left of the
threshold, and weighting by proximity to the threshold\medskip{}
\item The bandwidth $h$ determines how quickly weight falls off away from
the discontinuity\medskip{}
\end{itemize}
\end{frame}

\begin{frame}{Local Linear Regression}

\begin{center}
$\left(\hat{\alpha}_{0},\hat{\delta}_{0}\right)={\displaystyle \arg\min_{\alpha_{0},\delta_{0}}}{\displaystyle \sum_{i}}1\left\{ R_{i}\leq c\right\} K\left(\tfrac{R_{i}-c}{h}\right)\left(Y_{i}-\alpha_{0}-\delta_{0}\left(R_{i}-c\right)\right)^{2}$
\par\end{center}
\begin{itemize}
\item $\hat{\alpha}_{0}$ is an estimate of ${\displaystyle \lim_{r\rightarrow c^{-}}E\left[Y_{i}|R_{i}=r\right]}$\medskip{}
\item Run similar regression for the sample with $R_{i}>c$ to get $\hat{\alpha}_{1}$\medskip{}
\item Treatment effect estimate is $\hat{\alpha}_{1}-\hat{\alpha}_{0}$
\end{itemize}
\end{frame}

\begin{frame}{RD Implementation Issues}

\begin{itemize}
\item Remaining questions:\medskip{}

\begin{itemize}
\item Global or local?\medskip{}
\item How to choose order of polynomial?\medskip{}
\item How to choose bandwidth?\medskip{}
\end{itemize}
\item We will come back to these issues
\end{itemize}
\end{frame}

\begin{frame}{Fuzzy RD}

\begin{itemize}
\item In reality treatment is often generated by a discontinuous assignment
rule that isn't deterministic\medskip{}
\item Suppose that 
\end{itemize}
\begin{center}
${\displaystyle \lim_{r\rightarrow c^{-}}Pr\left[T_{i}=1|R_{i}=r\right]\neq\lim_{r\rightarrow c^{+}}Pr\left[T_{i}=1|R_{i}=r\right]}$
\par\end{center}
\begin{itemize}
\item The probability of treatment jumps at $R_{i}=c$, but not necessarily
from zero to one \medskip{}
\item This is a \textbf{fuzzy RD} scenario because treatment is only partly
determined by the threshold\medskip{}
\item Example: An income threshold determines eligibility for a government
program, but not every eligible household participates
\end{itemize}
\end{frame}

\setbeamercolor{background canvas}{bg=} \includepdf{aux/RDD2.pdf}
\begin{frame}{Fuzzy RD Assumptions}

\begin{itemize}
\item As before, assume the distributions of $Y_{i}(1)$ and $Y_{i}(0)$
are smooth around the threshold\medskip{}
\item Let $T_{i}(1)$ and $T_{i}(0)$ denote potential treatment statuses
for individual $i$ if s/he were located above and below the threshold.
Assume these are also smooth across the threshold, and
\end{itemize}
\begin{center}
$T_{i}(1)\geq T_{i}(0)\ \forall i$
\par\end{center}
\medskip
\begin{itemize}
\item Crossing the threshold weakly increases the probability of treatment
for everyone
\end{itemize}
\end{frame}

\begin{frame}{Fuzzy RD}

\begin{itemize}
\item Again compare individuals above and below the threshold:
\end{itemize}
\begin{center}
${\displaystyle \lim_{r\rightarrow c^{+}}E\left[Y_{i}|R_{i}=r\right]-{\displaystyle \lim_{r\rightarrow c^{-}}E\left[Y_{i}|R_{i}=r\right]}}$
\par\end{center}

\begin{center}
$={\displaystyle \lim_{r\rightarrow c^{+}}E\left[Y_{i}(0)+(Y_{i}(1)-Y_{i}(0))T_{i}|R_{i}=r\right]}$
\par\end{center}

\begin{center}
$-{\displaystyle \lim_{r\rightarrow c^{-}}E\left[Y_{i}(0)+(Y_{i}(1)-Y_{i}(0))T_{i}|R_{i}=r\right]}$
\par\end{center}

\begin{center}
{\small{}$={\displaystyle \lim_{r\rightarrow c^{+}}E\left[(Y_{i}(1)-Y_{i}(0))T_{i}(1)|R_{i}=r\right]-{\displaystyle \lim_{r\rightarrow c^{-}}E\left[(Y_{i}(1)-Y_{i}(0))T_{i}(0)|R_{i}=r\right]}}$}{\small\par}
\par\end{center}

\begin{center}
$=E\left[\left(Y_{i}(1)-Y_{i}(0)\right)\left(T_{i}(1)-T_{i}(0)\right)|R_{i}=c\right]$
\par\end{center}

\end{frame}

\begin{frame}{Fuzzy RD}

\begin{itemize}
\item Using monotonicity we can rewrite this as
\end{itemize}
\begin{center}
$=E\left[\left(Y_{i}(1)-Y_{i}(0)\right)\left(T_{i}(1)-T_{i}(0)\right)|R_{i}=c\right]$
\par\end{center}

\begin{center}
$=E\left[Y_{i}(1)-Y_{i}(0)|T_{i}(1)>T_{i}(0),R_{i}=c\right]\times Pr\left[T_{i}(1)>T_{i}(0)|R_{i}=c\right]$
\par\end{center}
\begin{itemize}
\item The compliance probability at the threshold is given by
\end{itemize}
\begin{center}
${\displaystyle \lim_{r\rightarrow c^{+}}E\left[T_{i}|R_{i}=r\right]-}{\displaystyle \lim_{r\rightarrow c^{-}}E\left[T_{i}|R_{i}=r\right]}=Pr\left[T_{i}(1)>T_{i}(0)|R_{i}=c\right]$
\par\end{center}

\end{frame}

\begin{frame}{Fuzzy RD}

\begin{itemize}
\item Then by taking a ratio we obtain
\end{itemize}
\begin{center}
$\dfrac{{\displaystyle \lim_{r\rightarrow c^{+}}E\left[Y_{i}|R_{i}=r\right]-{\displaystyle \lim_{r\rightarrow c^{-}}E\left[Y_{i}|R_{i}=r\right]}}}{{\displaystyle \lim_{r\rightarrow c^{+}}E\left[T_{i}|R_{i}=r\right]-{\displaystyle \lim_{r\rightarrow c^{-}}E\left[T_{i}|R_{i}=r\right]}}}=E\left[Y_{i}(1)-Y_{i}(0)|T_{i}(1)>T_{i}(0),R_{i}=c\right]$
\par\end{center}
\begin{itemize}
\item The ratio of the jump in the outcome CEF to the jump in the treatment
probability identifies an average treatment effect for individuals
who switch treatment status at the threshold\medskip{}
\item Look familiar?
\end{itemize}
\end{frame}

\begin{frame}{Fuzzy RD is IV}

\begin{itemize}
\item Fuzzy RD is IV using a threshold indicator $Z_{i}=1\left\{ R_{i}>c\right\} $
as an instrument for treatment in the neighborhood of the threshold
\end{itemize}
\end{frame}

\begin{frame}{RD is doubly LATE}

\begin{itemize}
\item This implies fuzzy RD estimates are local in two senses\medskip{}
\item First, they are local to the threshold, $R_{i}=c$\medskip{}

\begin{itemize}
\item Also applies to strict RD estimates\medskip{}
\end{itemize}
\item Second, they apply only to compliers at the threshold, rather than
everyone with $R_{i}=c$\medskip{}

\begin{itemize}
\item This is the ``local'' in LATE
\end{itemize}
\end{itemize}
\end{frame}

\begin{frame}{Fuzzy RD Implementation}

\begin{itemize}
\item Can implement fuzzy RD with a global polynomial 2SLS approach:
\end{itemize}
\begin{center}
$T_{i}=\lambda+\pi1\left\{ R_{i}>c\right\} +\sum_{k=1}^{K}\theta_{0k}(R_{i}-c)^{k}+\sum_{k=1}^{K}\theta_{1k}1\left\{ R_{i}>c\right\} \left(R_{i}-c\right)^{k}+\eta_{i}$
\par\end{center}

\begin{center}
$Y_{i}=\alpha+\beta\hat{T}_{i}+\sum_{k=1}^{K}\gamma_{0k}(R_{i}-c)^{k}+\sum_{k=1}^{K}\gamma_{1k}1\left\{ R_{i}>c\right\} (R_{i}-c)^{k}+\epsilon_{i}$
\par\end{center}
\begin{itemize}
\item Excluded instrument is $1\left\{ R_{i}>c\right\} $\medskip{}
\item Alternatively, can estimate each of the four limits in the Wald ratio
by local linear regressions of $Y_{i}$ and $T_{i}$ on $R_{i}$
\end{itemize}
\end{frame}

\begin{frame}{Visual RD}

\begin{itemize}
\item RD is appealing because the design is relatively simple and transparent\medskip{}
\item Key relationships are between the outcome and running variable, and
treatment and running variable\medskip{}
\item These relationships can be plotted, and one can look visually for
evidence of a discontinuity\medskip{}
\item In practice people won't believe an RD estimate if the result isn't
visible in a plot

\begin{itemize}
\item Automated plotting procedure: ``rdbinselect'' (Calonico, Cattaneo
and Titiunik, 2014)
\end{itemize}
\end{itemize}
\end{frame}
\begin{frame}{RD Diagnostics}

\begin{itemize}
\item Fundamental RD identifying assumption: potential outcome distributions
are smooth around the threshold\medskip{}
\item $1\left\{ R_{i}>c\right\} $ must be as good as randomly assigned
in the neighborhood of $R_{i}=c$\medskip{}
\item May be violated if individuals can exactly control the value of $R_{i}$
and therefore location relative to the threshold\medskip{}
\item Example: suppose income is self-reported, and people intentionally
misreport to gain program eligibility\medskip{}
\item Identifying assumption is untestable, but some common diagnostics
can serve as a guide to its plausibility
\end{itemize}
\end{frame}

\begin{frame}{RD Diagnostic I: Covariate Balance}

\begin{itemize}
\item If people are sorting in the neighborhood of the threshold, we'd expect
to see imbalance in distributions of observed covariates \medskip{}
\item This motivates a check for whether there is a discontinuity in $E\left[X_{i}|R_{i}=r\right]$
at $R_{i}=c$\medskip{}
\item Exactly analogous to a balance check in an RCT\medskip{}
\item Implement by using $X_{i}$ as the RD dependent variable
\end{itemize}
\end{frame}

\begin{frame}{RD Diagnostic II: Bunching}

\begin{itemize}
\item McCrary (2008) proposes a test for sorting that might compromise the
validity of an RD design\medskip{}
\item If individuals are strategically locating above or below the threshold,
we'd expect ``bunching'' on whichever side of the discontinuity
is preferable\medskip{}
\item More generally, strategic manipulation may generate anomalies in the
distribution of $R_{i}$ around the threshold\medskip{}
\item McCrary suggests looking for a discontinuity in the density of $R_{i}$
near $R_{i}=c$\medskip{}
\item This is often implemented visually\medskip{}
\item Bunching around the threshold compromises RD estimates of treatment
effects, but reflects a behavioral response that may be of substantive
interest, see Saez's research on bunching at kink point of tax schedule for instance.
\end{itemize}
\end{frame}
\begin{frame}
\frametitle{Urquiola and Verhoogen (2009)}
\setbeamercolor{background canvas}{bg=} \includepdf{aux/UV_f1.pdf}
\end{frame}

\setbeamercolor{background canvas}{bg=} \includepdf{aux/UV_f2.pdf}

\setbeamercolor{background canvas}{bg=} \includepdf{aux/UV_f3.pdf}

\setbeamercolor{background canvas}{bg=} \includepdf{aux/UV_f4.pdf}

\begin{frame}{RD Estimation: Global vs. Local}

\begin{itemize}
\item Let's return to our discussion of RD implementation\medskip{}
\item When should we use a global polynomial approach vs. a nonparametric
local approach?\medskip{}
\item There is no real conceptual distinction: we need to choose a bandwidth
and a polynomial order\medskip{}
\item Global estimators use an infinite bandwidth and a relatively high-order
polynomial\medskip{}
\item Local estimators use a smaller bandwidth and a lower-order polynomial
(often linear)
\end{itemize}
\end{frame}

\begin{frame}{RD Estimation: Global vs. Local}

\begin{itemize}
\item Global estimators can rely heavily on data far from the discontinuity
(Gelman and Imbens 2014, ``Why high-order polynomials should not
be used in regression discontinuity designs'')\medskip{}
\item If the outcome CEF is well-approximated by a polynomial, then using
points far from the discontinuity to estimate the polynomial is legitimate
and increases precision. If not, the global approach may perform poorly\medskip{}
\item Not much is known about choosing the optimal order of a global polynomial
for estimating the causal effect of interest; in contrast, there is
a large literature on optimal bandwidths for local estimators of treatment
effects\medskip{}
\item Most researchers have gravitated towards the local approach
\end{itemize}
\end{frame}

\begin{frame}{Kernels and Bandwidths}

\begin{itemize}
\item Suppose we want to run local linear regression. What kernel and bandwidth
should we choose?\medskip{}
\item Many researchers use the ``edge'' kernel:
\end{itemize}
\begin{center}
$K(u)=1\left\{ |u|\leq1\right\} \cdot(1-|u|)$
\par\end{center}
\begin{itemize}
\item The edge kernel has optimality properties in boundary estimation problems
(Cheng et al., 1997)\medskip{}
\item Also intuitively appealing: generates weighting function
\end{itemize}
\begin{center}
$K\left(\tfrac{R_{i}-c}{h}\right)=1\left\{ |R_{i}-c|\leq h\right\} \cdot\left(1-\tfrac{|R_{i}-c|}{h}\right)$
\par\end{center}
\begin{itemize}
\item The bandwidth $h$ can then be interpreted as cutoff distance beyond
which data are not used, and weights fall linearly from 1 to 0 in
remaining sample
\end{itemize}
\end{frame}

\begin{frame}{Optimal Bandwidth}

\begin{itemize}
\item Recent literature focuses on optimal choice of bandwidth $h$\medskip{}
\item Bias/variance tradeoff: Smaller bandwidth reduces bias from using
points away from the boundary, but increases precision\medskip{}
\item Intuitively, if there is not a lot of curvature in the CEF of $Y_{i}$,
the bias from using points away from the boundary to estimate a regression
will be small\medskip{}
\item Imbens and Kalyanaraman (IK, 2012) use an asymptotic approximation
to the mean squared error of the RD estimator and derive the MSE-minimizing
bandwidth\medskip{}
\item The optimal bandwidth depends on the curvature of the CEF near the
discontinuity -- IK propose to use plug-in estimators of parameters
governing curvature
\end{itemize}
\end{frame}

\begin{frame}{Robust Confidence Intervals}

\begin{itemize}
\item The IK bandwidth minimizes MSE and is therefore well-suited to estimation
\medskip{}
\item Calonico, Cattaneo and Titiunik (CCT, 2014) show that it is poorly
suited for inference, however: the IK bandwidth leaves an asymptotically
non-negligible bias term in the estimate, so naive inference
can lead to misleading confidence intervals \medskip{}
\item CCT advocate using a second, smaller bandwidth that removes this bias
term when constructing confidence intervals
\end{itemize}
\end{frame}

\begin{frame}{Robust Confidence Intervals}

\begin{itemize}
\item Unfortunately, the CCT bandwidth is typically much smaller than the
IK bandwidth and generates large confidence intervals\medskip{}
\item The optimal bandwidth literature may limit the feasibility of RD projects
in practice since lots of data is required to detect an effect with
CCT's robust confidence intervals\medskip{}
\item The IK bandwidth, CCT confidence intervals, and corresponding estimators
for fuzzy RD designs are automated in the \emph{rdrobust} stata package
\end{itemize}
\end{frame}
%%
\begin{frame}[plain]{Is this a compelling RD-plot?}
\begin{center}
\includegraphics[scale=0.7]{aux/kirabobetter1}
\par\end{center}
\end{frame}
%%
\begin{frame}[plain]{Can you guys guess what is the t-stat in this graph?}
\begin{center}
\includegraphics[scale=0.7]{aux/kirabobetter}
\par\end{center}
\end{frame}
%%

%%
%%
%%
\begin{frame}[plain]{RD with a Discrete Running Variable}
\pause
\begin{center}
\includegraphics[height=0.82\textheight,keepaspectratio]{aux/card_lee.jpg}
\par\end{center}
\end{frame}
%%

\begin{frame}{RD with a Discrete Running Variable}
\begin{itemize}
\item In some cases the running variable is discrete rather than continuous\medskip{}
\item This can occur either because the true running variable is continuous
but the data are grouped, or because assignment occurs at the group
level\medskip{}
\item Example: Medicare eligibility begins at age 65, but age is typically
recorded in months or years\medskip{}
\item With data of this type it is not possible to compare outcomes in arbitrarily
small neighborhoods of the cutoff; identification requires some form
of parametric extrapolation\medskip{}
\item Card and Lee (2008): recommend clustering on values of the running variable in settings
with a group-level running variable
\end{itemize}
\end{frame}
%%
\begin{frame}[plain]{Sitting on the frontier}
\begin{center}
\includegraphics[scale=0.3]{aux/kolesar}
\par\end{center}
\end{frame}
\begin{frame}[plain]{Sitting on the frontier}
\begin{center}
\includegraphics[scale=0.3]{aux/honest.jpg}
\par\end{center}
\end{frame}
%%
\begin{frame}{RD Example: Clark and Martorell (2014)}

\begin{itemize}
\item Clark and Martorell (JPE 2014) use an RD design to estimate the causal
effect of high school graduation on earnings \medskip{}
\item Two views on the causal effect of schooling on earnings:\medskip{}

\begin{itemize}
\item Human capital: Schooling raises productivity\medskip{}
\item Signaling: Schooling reveals ability but has no productive value\medskip{}
\end{itemize}
\item OLS returns to education are especially large for grade 12\medskip{}
\item How much of this ``sheepskin effect'' reflects signaling?
\end{itemize}
\end{frame}

\begin{frame}{Clark and Martorell (2014)}

\begin{itemize}
\item CM use the fact that students in Texas must pass exams before graduating
high school\medskip{}
\item Testing starts in 10th grade and students can try multiple times,
but eventually face a ``last chance'' exam at the end of 12th grade\medskip{}
\item Students who just barely fail vs. barely pass should have similar
human capital, but differ in educational credentials \medskip{}
\item RD therefore plausibly identifies the signaling value of a diploma\medskip{}
\item There is some ``slippage'' even with last-chance exams -- so the
RD is fuzzy
\end{itemize}
\end{frame}

\setbeamercolor{background canvas}{bg=} \includepdf{aux/mc_f1.pdf}

\setbeamercolor{background canvas}{bg=} \includepdf{aux/mc_t2.pdf}

\setbeamercolor{background canvas}{bg=} \includepdf{aux/mc_f2aa.pdf}

\setbeamercolor{background canvas}{bg=} \includepdf{aux/mc_f2a.pdf}

\setbeamercolor{background canvas}{bg=} \includepdf{aux/mc_f2b.pdf}

\setbeamercolor{background canvas}{bg=} \includepdf{aux/mc_f2c.pdf}

\setbeamercolor{background canvas}{bg=} \includepdf{aux/mc_t4.pdf}

\setbeamercolor{background canvas}{bg=} \includepdf{aux/mc_t5.pdf}
\begin{frame}
\centering
\LARGE{Encore}
\end{frame}

\begin{frame}[plain]{Multiple Discontinuities}

\begin{itemize}
\item Often more than one discontinuity present. 
\end{itemize}
\medskip{}

\begin{itemize}
\item Nonparametric solution is to estimate each RD separately then average
across them as desired. (see Black, Galdo, Smith 2007 AER) 

\begin{itemize}
\item Only feasible if discontinuities far enough apart to get a good estimate
on either side 
\end{itemize}
\end{itemize}
\medskip{}

\begin{itemize}
\item Second approach is to impose some sort of stationarity -- e.g. running
variable becomes distance to nearest discontinuity. 
\end{itemize}
\medskip{}

\begin{itemize}
\item Third approach is to go fully parametric.
\end{itemize}
 See Angrist and Lavy (1999) for an example.
\end{frame}


\begin{frame}{Regression Kink Design}

\begin{itemize}
\item Recent extension of RD: The\textbf{ regression kink design }(RKD;
Card et al., 2015) \medskip{}
\item Instead of exploiting a discontinuity in the CEF of the treatment
variable, the regression kink design exploits a kink in the CEF of
a continuous treatment (i.e. a discontinuity in the first derivative)\medskip{}
\item A corresponding kink in the distribution of the outcome variable suggests
the presence of a treatment effect
\end{itemize}
\end{frame}

\begin{frame}{Regression Kink Design}

\begin{itemize}
\item Suppose the treatment of interest is a deterministic function of the
running variable:
\end{itemize}
\begin{center}
$S_{i}=b\left(R_{i}\right)$
\par\end{center}
\begin{itemize}
\item Here $b(\cdot)$ is a continuous function with a kink at $c$\medskip{}
\item Let $f_{i}(s)$ denote $i$'s potential outcome as a function of the
treatment. The observed outcome is
\end{itemize}
\begin{center}
$Y_{i}=f_{i}(S_{i})$
\par\end{center}

\end{frame}

\begin{frame}{Regression Kink Design}

\begin{itemize}
\item Then under mild regularity conditions:
\end{itemize}
\begin{center}
$\dfrac{{\displaystyle \lim_{r\rightarrow c^{+}}\tfrac{dE\left[Y_{i}|R_{i}=r\right]}{dr}-\lim_{r\rightarrow c^{-}}\tfrac{dE\left[Y_{i}|R_{i}=r\right]}{dr}}}{{\displaystyle \lim_{r\rightarrow c^{+}}b^{\prime}(r)}-{\displaystyle \lim_{r\rightarrow c^{-}}}b^{\prime}(r)}=E\left[f_{i}^{\prime}(S_{i})|R_{i}=c\right]$
\par\end{center}
\begin{itemize}
\item The ratio of the discontinuity in the outcome derivative to the discontinuity
in the treatment derivative identifies the average marginal effect
of treatment for individuals at the threshold\medskip{}
\item As before, the key assumption is that potential outcomes are smooth
around the threshold -- any kink in the outcome CEF must be due to
the treatment\medskip{}
\item Diagnostics: Look for kinks in covariate distributions, or bunching
in the density of $R_{i}$
\end{itemize}
\end{frame}

\begin{frame}{Fuzzy Regression Kink Design}

\begin{itemize}
\item As with RD, we can generalize RKD to a ``fuzzy'' scenario where
the treatment is not a deterministic function of $R_{i}$
\end{itemize}
\begin{center}
$S_{i}=b_{i}(R_{i})$
\par\end{center}
\begin{itemize}
\item $b_{i}(r)$ may not depend on $r$ for some individuals, and may have
a continuous derivative for others
\item Monotonicity condition:
\end{itemize}
\begin{center}
${\displaystyle \lim_{r\rightarrow c^{+}}b_{i}^{\prime}(r)\geq{\displaystyle \lim_{r\rightarrow c^{-}}b_{i}^{\prime}(r)\ \forall i}}$
\par\end{center}
\begin{itemize}
\item All kinks in the treatment assignment rule work in the same direction
\end{itemize}
\end{frame}

\begin{frame}{Fuzzy Regression Kink Design}

\begin{itemize}
\item Then under appropriate regularity conditions:
\end{itemize}
\begin{center}
$\dfrac{{\displaystyle \lim_{r\rightarrow c^{+}}\tfrac{dE\left[Y_{i}|R_{i}=r\right]}{dr}-\lim_{r\rightarrow c^{-}}\tfrac{dE\left[Y_{i}|R_{i}=r\right]}{dr}}}{{\displaystyle \lim_{r\rightarrow c^{+}}\tfrac{dE\left[S_{i}|R_{i}=r\right]}{dr}-\lim_{r\rightarrow c^{-}}\tfrac{dE\left[S_{i}|R_{i}=r\right]}{dr}}}=E\left[\omega_{i}\cdot f_{i}^{\prime}(S_{i})|R_{i}=c\right]$
\par\end{center}
\begin{itemize}
\item $\omega_{i}$ can be thought of as a ``complier'' weight that captures
the extent to which $i$ responds to the kink\medskip{}
\item Estimation issues with RKD parallel those with RD\medskip{}
\item Can be implemented via local polynomial regression with the analogue
of the IK bandwidth and CCT robust CI\medskip{}
\item Data requirements for nonparametric RKD estimation are substantial
\end{itemize}
\end{frame}

\begin{frame}{RKD Estimation}

\begin{itemize}
\item Estimation issues with RKD parallel those with RD

\begin{itemize}
\item Can be implemented via local polynomial regression with the analogues
of the IK bandwidth and CCT robust CI
\item \emph{rdrobust} provides an implementation \medskip{}
\end{itemize}
\item Corresponding diagnostics

\begin{itemize}
\item Visualization
\item Kinks in baseline covariates?
\end{itemize}
\medskip{}

\item Need even more data to precisely estimate kinks than to estimate discontinuities
\end{itemize}
\end{frame}

\end{document}
